\section{What is and is not proved about \texorpdfstring{$A(p)$}{A(p)}}
\label{a:sec:scope}

``Closed form'' must be made precise before any claim is stated, because at least
four inequivalent readings are in play and a result for one is not a result for
another:
\begin{enumerate}[leftmargin=*]
\item[\textbf{(V)}] the \emph{value} $A(p_0)$ is elementary, or algebraic, for a
      fixed rational $p_0$;
\item[\textbf{(E)}] the \emph{map} $p\mapsto A(p)$ is an elementary function of
      $p$;
\item[\textbf{(D)}] the map is \emph{differentially algebraic in $p$}, that is, it
      satisfies some algebraic ordinary differential equation in $p$;
\item[\textbf{(S)}] the map is expressible through standard special functions
      ($\Gamma$, $\zeta$, ${}_2F_1$, theta or $q$-series).
\end{enumerate}
\Cref{a:thm:ht} speaks to the function $z\mapsto\psi_r(z)$ at fixed $r$, and
licenses exactly this much: \emph{no evaluation of $A(p)$ can be short-circuited
through an algebraic differential equation for the lineariser, at any subcritical
parameter}. It does not, by itself, decide any of (V), (E), (D) or (S). Two gaps
must be stated plainly.

\subsection*{The value gap}

A hypertranscendental function can take elementary --- even algebraic --- values at
individual points: $\Gamma$ is hypertranscendental by Hölder's theorem, and yet
$\Gamma(\tfrac12)=\sqrt\pi$. Nothing established so far forbids, say, $A(\tfrac14)$
from being an algebraic number.

There \emph{is} a genuine values theorem in this circle of ideas. Becker and
Bergweiler proved in 1993 \cite{Becker1993} that the Böttcher conjugacy between two
polynomials of the \emph{same degree} $d\ge2$ with algebraic coefficients, not
linearly conjugate to monomials or Chebyshev polynomials, is transcendental and
takes transcendental values at algebraic points; Mahler's method drives the proof
(for the method see Nishioka \cite{Nishioka1996}). One might hope to transfer this
to the Koenigs value $\psi_r(\tfrac12)$. It does not transfer, for two independent
reasons.

\emph{Formally}, the Schröder equation pairs an inner map of degree two with an
outer map --- multiplication by $r$ --- of degree one, so the equal-degree
hypothesis fails outright.

\emph{Structurally}, no easy repair exists. Mahler's method balances
doubly-exponential analytic decay along an orbit against comparable height growth,
and in the Böttcher setting both exponents run like $d^{\,n}$. In the attracting
Koenigs setting the orbit $w_n = f_r^{\circ n}(\tfrac12)$ decays only geometrically,
$|w_n|\sim A\,r^{n}$, while the heights of the rational numbers $w_n$ still grow
doubly exponentially, $h(w_n)\sim C\,2^{\,n}$; the resulting inequalities run the
wrong way at every $n$. Note that the obstruction is \emph{not} the germ/basin
issue of \cref{a:rem:basin} --- the pullback there moves the evaluation point
arbitrarily deep into the germ at the cost of a rational factor --- it is the degree
structure of the equation itself.

The consequence deserves emphasis:
\begin{center}
\emph{transcendence --- indeed, even irrationality --- of $A(p_0)$ is at present
not proved\\ for a single rational $p_0\in(0,\tfrac12)$.}
\end{center}

\subsection*{The parameter gap}

\Cref{a:thm:ht} says nothing about the map $p\mapsto A(p)$: readings (E) and (D)
are open. To the best of our knowledge the question is not merely unsolved but
unaddressed in the literature. The natural machinery is the parametrised
differential Galois theory of difference equations of Hardouin and Singer
\cite{HardouinSinger2008}, designed precisely to detect differential relations in
auxiliary directions for linear difference equations such as $\sigma\psi = r\psi$;
it does not apply as it stands, because the derivation $\partial_r$ fails to commute
with the substitution operator $\sigma\colon w\mapsto f_r(w)$, which itself moves
with $r$. Nor does any feature of the near-critical asymptotics obstruct (E) or
(D): the logarithmic corrections to \cref{a:eq:Anearcritical} are themselves
elementary functions of $\varepsilon$.

\section{Inverse-symbolic evidence at eleven rational parameters}
\label{a:sec:pslq}

Since the theorems leave (V) open, the value question was attacked directly, by
integer-relation detection. The values
\begin{equation}
\label{a:eq:pslqparams}
A(p),\qquad
p\in\Bigl\{\tfrac1{10},\tfrac18,\tfrac16,\tfrac15,\tfrac14,\tfrac3{10},
\tfrac13,\tfrac38,\tfrac25,\tfrac5{12},\tfrac9{20}\Bigr\},
\end{equation}
were computed to more than $1150$ significant digits by two mechanically
independent routes --- the exact product identity \cref{a:eq:prodFormula} iterated
to convergence, and a pullback of the orbit into the Koenigs germ followed by
evaluation of the Koenigs series --- which agreed to at least $1192$ digits at every
parameter. Each value was then subjected to a PSLQ battery at $\ge200$-digit
effective rigour:
\begin{itemize}[leftmargin=*]
\item \emph{algebraicity}: no integer polynomial relation of degree $\le10$ and
      coefficient height $\le10^{6}$;
\item \emph{linearity}: no rational linear relation, at height $\le10^{8}$, against
      the constants $\pi$, $\pi^{2}$, $\ln2$, $\ln3$, $\ln5$, $\gamma$, $\zeta(3)$,
      $\sqrt2$, $\sqrt3$, $\sqrt5$, $\ee$, $\ee^{\pi}$, $\Gamma(\tfrac13)$,
      $\Gamma(\tfrac14)$ and Catalan's constant --- also tested for $1/A$, the
      quasi-stationary mean itself, and for $A/\varepsilon$;
\item \emph{bilinearity}: no representation $A = L_1/L_2$ with $L_1,L_2$ integer
      linear forms of height $\le10^{8}$ in a twelve-constant basis, run at $460$
      digits;
\item \emph{multiplicativity}: no relation
      $A = \ee^{q_0}2^{q_1}3^{q_2}5^{q_3}7^{q_4}\pi^{q_5}$ with rational exponents
      of height $\le10^{10}$;
\item \emph{cross-parameter}: no rational linear or multiplicative relation linking
      the eleven values to one another, run at $500$ digits.
\end{itemize}
In all, $123$ tests returned no relation --- indeed no discovery-stage candidate:
every PSLQ run terminated on its internal norm bound, so each null carries a
genuine ``no relation up to the stated height'' guarantee. Throughout, the
precision head-room, digits divided by basis size, exceeded every height cap by at
least eight orders of magnitude, so a returned relation would have been meaningful
rather than numerical noise.\footnote{Computations in \texttt{mpmath} at working
precisions of $310$--$1200$ digits; the values, scripts and full per-test log are
archived with the project sources and are not reproduced in this thesis.}

This is finite evidence, and its scope must be stated honestly: it excludes
low-complexity closed forms over the tested constants, and nothing more. A closed
form over untested constants --- $\Gamma(\tfrac17)$, say, or periods of
higher-genus curves --- or one of height above the caps would have escaped the
search. It is evidence for the conjecture, not a proof of it.

\subsection{The working claim}
\label{a:sec:workingclaim}

\begin{quote}
For every fixed $r=2p\in(0,1)$ the Koenigs linearisation $\psi_r$ is
hypertranscendental in its own variable (\cref{a:thm:ht}) --- a theorem, with no
exceptional parameters in the attracting window, since differentially algebraic
Schröder functions exist only over repelling fixed points. Consequently no
elementary expression in $z$ for the conjugacy can be evaluated to produce $A(p)$.
Whether individual values $A(p_0)$ are transcendental, or even irrational, and
whether the map $p\mapsto A(p)$ is elementary or differentially algebraic in $p$,
remain open questions. No closed form under any reading is known, none is expected,
and eleven rational values resist inverse-symbolic identification against standard
constants at over $200$-digit rigour.
\end{quote}

Full derivations of the elementary solutions at $r=2$ and $r=4$, and their place
inside the Becker--Bergweiler classification, are collected in
\cref{a:app:integrable}.
